% ============================================================================
%  WO-CHART-FIELD-B --- Unified Scientific Data Manifold
%  VSEPR-SIM Theoretical Division
%  Classification: TOP SECRET (internal development document)
%
%  Subject: Canonical mathematical specification of the unified scientific
%           data manifold underpinning .chart, .xchart, .dynx, .flux, and
%           .x artifact types; domain manifolds for thermodynamics, MD,
%           flux fields, radiation transport, excitation dynamics, and the
%           effective gluon/color-field proxy layer; Qt desktop stack;
%           and the universal validation and implementation contract.
%
%  Compile: pdflatex (twice for ToC)
% ============================================================================
\documentclass[11pt, a4paper, oneside]{article}

\usepackage[top=1.0in, bottom=1.0in, left=1.15in, right=1.15in,
            headheight=14pt]{geometry}
\usepackage[utf8]{inputenc}
\usepackage[T1]{fontenc}
\usepackage{lmodern}
\usepackage{microtype}
\usepackage{amsmath, amssymb, amsthm, mathtools}
\usepackage{bm}
\usepackage{booktabs}
\usepackage{array}
\usepackage{tabularx}
\usepackage{listings}
\usepackage{xcolor}
\usepackage{enumitem}
\usepackage{fancyhdr}
\usepackage{titlesec}
\usepackage[colorlinks=true, linkcolor=black, citecolor=black,
            urlcolor=blue]{hyperref}

% --- Code style --------------------------------------------------------------
\definecolor{codegray}{rgb}{0.5,0.5,0.5}
\definecolor{codeblue}{rgb}{0.13,0.13,0.9}
\definecolor{codegreen}{rgb}{0.0,0.5,0.0}
\definecolor{codebg}{rgb}{0.97,0.97,0.97}
\lstset{
  backgroundcolor=\color{codebg},
  basicstyle=\ttfamily\small,
  keywordstyle=\color{codeblue}\bfseries,
  commentstyle=\color{codegreen}\itshape,
  numberstyle=\tiny\color{codegray},
  numbers=left, numbersep=6pt,
  breaklines=true, frame=single,
  rulecolor=\color{codegray},
  captionpos=b, showstringspaces=false,
  tabsize=2
}

% --- Headers -----------------------------------------------------------------
\pagestyle{fancy}
\fancyhf{}
\fancyhead[L]{\small\textit{WO-CHART-FIELD-B --- Unified Scientific Data Manifold}}
\fancyhead[R]{\small\textit{VSEPR-SIM / TOP SECRET}}
\fancyfoot[C]{\thepage}
\renewcommand{\headrulewidth}{0.4pt}

% --- Section style -----------------------------------------------------------
\titleformat{\section}{\Large\bfseries}{\S\thesection}{1em}{}
\titleformat{\subsection}{\large\bfseries}{\thesubsection}{1em}{}
\titleformat{\subsubsection}{\normalsize\bfseries}{\thesubsubsection}{1em}{}

% --- Math convenience --------------------------------------------------------
\newcommand{\ang}{\text{\AA}}
\newcommand{\fs}{\text{fs}}
\newcommand{\eV}{\text{eV}}
\newcommand{\amu}{\text{amu}}
\newcommand{\Rec}{\mathcal{R}}           % universal record
\newcommand{\Tab}{\mathcal{T}}           % table of records
\newcommand{\Proj}{\Pi}                  % projection operator
\newcommand{\Val}{\mathcal{V}}           % validation operator
\newcommand{\Bund}{\mathcal{B}}          % .x bundle
\newcommand{\Qvec}{\mathbf{Q}}           % universal state vector
\newcommand{\xvec}{\mathbf{x}}
\newcommand{\vvec}{\mathbf{v}}
\newcommand{\Fvec}{\mathbf{F}}
\newcommand{\Phivec}{\bm{\Phi}}
\newcommand{\uvec}{\mathbf{u}}
\newcommand{\phibold}{\bm{\phi}}
\newcommand{\Hvec}{\mathbf{H}}           % lattice matrix (context-dep.)
\newcommand{\avec}{\mathbf{a}}
\newcommand{\Omvec}{\bm{\Omega}}         % radiation direction
\DeclareMathOperator{\Hash}{Hash}
\newcommand{\norm}[1]{\left\|#1\right\|}

% ============================================================================
\begin{document}

\begin{titlepage}
  \centering
  \vspace*{1.2cm}
  {\large\scshape VSEPR-SIM Theoretical Division\par}
  \vspace{0.3cm}
  {\small TOP SECRET --- Internal Development Document\par}
  \vspace{1.2cm}
  {\Huge\bfseries WO-CHART-FIELD-B\par}
  \vspace{0.4cm}
  {\LARGE Unified Scientific Data Manifold\par}
  \vspace{0.8cm}
  {\large\itshape
    Chart \& XChart artifacts\,|\,Dynamic event streams\,|\,\\[0.3em]
    Flux fields\,|\,Radiation transport\,|\,Excitation dynamics\,|\,\\[0.3em]
    Effective gluon proxy\,|\,Qt desktop stack\,|\,\\[0.3em]
    Validation mathematics\,|\,Universal implementation contract\par}
  \vspace{1.4cm}
  {\normalsize VSEPR-SIM Development Work Order Series\par}
\end{titlepage}

\tableofcontents
\newpage

% ============================================================================
\section{Core Thesis and Motivation}
\label{sec:thesis}
% ============================================================================

\subsection{The plotting-scale artifact problem}

The previous Python plotting layer produced outputs such as a heat capacity
curve for the hydrogen atom with a y-axis offset of the form:

\[
  \underbrace{1 \times 10^{-10}}_{\text{numerical noise}} + 2.078603 \times 10^{1}
\]

The axis is communicating that $C_p \approx 20.786\ \mathrm{J\,mol^{-1}\,K^{-1}}$
(the correct monatomic ideal-gas value $\frac{5}{2}R$), but the floating-point
residuals in the time-integrated trajectory are being treated as physically
meaningful variation and plotted on a collapsed scale. The physics is correct.
The representation is not.

This is not a cosmetic problem. A plot that suppresses its own axis scale or
offsets it to highlight noise carries no reproducible scientific information.
It is missing metadata about:

\begin{itemize}
  \item what equation and coefficient set produced the curve,
  \item what phase assumption and species model were used,
  \item what unit convention the y-axis values are in,
  \item what validity range the model applies to,
  \item whether the axis offset was intentional or a numerical artifact.
\end{itemize}

A PNG image cannot answer any of these questions. The solution is to make
the chart itself a structured scientific object that carries its own
provenance, validation state, and projection definition alongside the data.

\subsection{The unified abstraction}

Every output object in VSEPR-SIM is a sampled field over a generalised domain.
Steam tables, radiation transport fields, MD trajectories, flux surfaces,
excitation timelines, and effective gluon event streams are all projections of
the same universal state vector $\Qvec$. One projection operator, all output
types. The same validation contract applies to all of them.

\medskip
\noindent\textbf{File type stack affected by this work order:}

\begin{center}
\begin{tabular}{ll}
  \toprule
  Type & Role \\
  \midrule
  \texttt{.chart}  & Lightweight plot recipe \\
  \texttt{.xchart} & Rich chart + data + provenance visualization object \\
  \texttt{.dynx}   & Dynamic force, bond, vector, excitation, event stream \\
  \texttt{.flux}   & Scalar/vector/tensor field state \\
  \texttt{.x}      & Universal session / artifact bundle \\
  \bottomrule
\end{tabular}
\end{center}

\noindent\textbf{Work order parts:}

\begin{description}[leftmargin=2em]
  \item[Part A] Add \texttt{.chart} and \texttt{.xchart} as first-class artifact formats.
  \item[Part B] Add Qt desktop support: 2D charts, 3D surfaces, vector fields,
    scalar field slices, radiation maps, excitation timelines, event trees,
    batch comparison dashboards.
  \item[Part C] Connect chart objects to: steam table presolved MD outputs,
    \texttt{.dynx} dynamic event streams, \texttt{.flux} field samples,
    \texttt{.x} universal session bundles.
  \item[Part D] Add validation rules preventing physically meaningless charts
    from being treated as verified outputs.
\end{description}

% ============================================================================
\section{Universal State Vector and Sample Record}
\label{sec:universal}
% ============================================================================

\subsection{Generalised domain}

Every sampled object lives on a product domain:
\begin{equation}
  \mathcal{D}
  = \Omega_x \times \Omega_t \times \Omega_\theta
    \times \Omega_E \times \Omega_s
  \label{eq:domain}
\end{equation}
where
\[
  \Omega_x \subseteq \mathbb{R}^3,\quad
  \Omega_t \subseteq \mathbb{R}_{\geq 0},\quad
  \Omega_\theta \subseteq \mathbb{R}^{n_\theta},\quad
  \Omega_E \subseteq \mathbb{R},\quad
  \Omega_s \subseteq \mathbb{R}^{n_s}.
\]

\subsection{Universal state vector}

The universal sampled state is:
\begin{equation}
  \Qvec = \bigl[\xvec,\; t,\; \bm{\theta},\; E,\; \mathbf{s},\;
                \vvec,\; \Fvec,\; \Phivec,\; \mathbf{M},\; \mathbf{H}\bigr]
  \label{eq:Qvec}
\end{equation}
with components:
\begin{align*}
  \xvec &= (x,y,z) & &\text{spatial position} \\
  t     & & &\text{time coordinate} \\
  \bm{\theta} &= (T, P, \rho, \ldots) & &\text{thermodynamic / control variables} \\
  E     & & &\text{energy coordinate} \\
  \mathbf{s} & & &\text{internal state vector} \\
  \vvec &= (v_x, v_y, v_z) & &\text{velocity} \\
  \Fvec &= (F_x, F_y, F_z) & &\text{force} \\
  \Phivec &= (\Phi_0, \Phi_x, \Phi_y, \Phi_z, \Phi_{xx}, \ldots) & &\text{field values} \\
  \mathbf{M} & & &\text{model metadata} \\
  \mathbf{H} & & &\text{provenance hash}
\end{align*}

\subsection{Primitive sampled record}

The primitive record for sample $i$ is:
\begin{equation}
  \Rec_i = \bigl(\mathrm{id}_i,\; \xvec_i,\; t_i,\; \bm{\theta}_i,\; E_i,\;
                  \mathbf{s}_i,\; \vvec_i,\; \Fvec_i,\; \Phivec_i,\;
                  w_i,\; \mu_i,\; H_i\bigr)
  \label{eq:record}
\end{equation}
where $w_i$ is the statistical weight and $\mu_i$ is the uncertainty or
variance metric.

\subsection{General table form}

The full dataset is the table of $N$ records:
\begin{equation}
  \Tab = \{\Rec_i\}_{i=1}^{N}
  \label{eq:table}
\end{equation}

In column form, $\Tab$ is the matrix:
\begin{equation}
  \Tab =
  \bigl[\;
    \mathrm{id}\;\Big|\;
    x\; y\; z\;\Big|\;
    t\;\Big|\;
    \theta_1 \cdots \theta_{n_\theta}\;\Big|\;
    E\;\Big|\;
    s_1 \cdots s_{n_s}\;\Big|\;
    v_x\; v_y\; v_z\;\Big|\;
    F_x\; F_y\; F_z\;\Big|\;
    \Phi_0 \cdots\;\Big|\;
    w\;\Big|\; \mu\;\Big|\; H
  \;\bigr]
  \label{eq:table_col}
\end{equation}

This is the mathematical basis of \texttt{.xchart}.

% ============================================================================
\section{Universal Projection Operator}
\label{sec:projection}
% ============================================================================

A chart is a projection of the table $\Tab$ into a renderable space:
\begin{equation}
  \Proj : \Tab \to \mathbb{R}^k
  \label{eq:proj}
\end{equation}

\subsection{2D chart}

\begin{equation}
  C_{\mathrm{2D}} = \bigl\{(a_i, b_i)\bigr\}_{i=1}^{N},
  \qquad
  a_i = f_a(\Rec_i),\quad b_i = f_b(\Rec_i)
  \label{eq:proj2d}
\end{equation}

\subsection{3D surface}

\begin{equation}
  C_{\mathrm{3D}} = \bigl\{(a_i, b_i, c_i)\bigr\}_{i=1}^{N},
  \qquad
  a_i = f_a(\Rec_i),\quad b_i = f_b(\Rec_i),\quad c_i = f_c(\Rec_i)
  \label{eq:proj3d}
\end{equation}

\subsection{Vector-field plot}

\begin{equation}
  C_V = \bigl\{(\xvec_i, \uvec_i)\bigr\}_{i=1}^{N},
  \qquad
  \uvec_i = f_u(\Rec_i)
  \label{eq:projvec}
\end{equation}

Typical choices: $\uvec_i = \vvec_i$ (velocity field), $\uvec_i = \Fvec_i$
(force field), $\uvec_i = \Phivec_{\mathrm{vec},i}$ (flux vector field).

\medskip
\noindent The reason not to build separate plotting systems for steam tables,
radiation fields, and gluon events is that the \emph{plot} is just a
projection. The data changes. The math does not.

% ============================================================================
\section{Steam Table / Thermodynamic Manifold}
\label{sec:steam}
% ============================================================================

\subsection{State vector}

\begin{equation}
  \mathbf{S}_{\mathrm{steam}}
  = \bigl[T,\; P,\; \rho,\; u,\; h,\; s,\; C_p,\; C_v,\;
           \mu_{\mathrm{visc}},\; k,\; x_q,\; \phi\bigr]
  \label{eq:steam_state}
\end{equation}

\begin{center}
\begin{tabular}{lll}
  \toprule
  Symbol & Quantity & Unit \\
  \midrule
  $T$             & Temperature           & K \\
  $P$             & Pressure              & Pa \\
  $\rho$          & Density               & kg/m$^3$ \\
  $u$             & Specific internal energy & J/kg \\
  $h$             & Specific enthalpy     & J/kg \\
  $s$             & Specific entropy      & J/(kg\,K) \\
  $C_p$           & Isobaric heat capacity & J/(mol\,K) \\
  $C_v$           & Isochoric heat capacity & J/(mol\,K) \\
  $\mu_{\mathrm{visc}}$ & Dynamic viscosity & Pa\,s \\
  $k$             & Thermal conductivity  & W/(m\,K) \\
  $x_q$           & Vapour quality        & --- \\
  $\phi$          & Phase code            & integer \\
  \bottomrule
\end{tabular}
\end{center}

\subsection{Governing property functions}

\begin{align}
  \rho &= \rho(T,P) & h &= h(T,P) \notag \\
  u    &= u(T,P)    & s &= s(T,P) \notag \\
  C_p  &= C_p(T,P)  & C_v &= C_v(T,P)
\end{align}

\subsection{Thermodynamic identities}

Enthalpy--internal energy relation:
\begin{equation}
  h = u + \frac{P}{\rho}
  \label{eq:enthalpy}
\end{equation}

Heat capacities via central finite differences over the MD-derived table:
\begin{align}
  C_p(T_i, P_j) &\approx \frac{h(T_{i+1}, P_j) - h(T_{i-1}, P_j)}{T_{i+1} - T_{i-1}}
  \label{eq:Cp_fd} \\[4pt]
  C_v(T_i, \rho_j) &\approx \frac{u(T_{i+1}, \rho_j) - u(T_{i-1}, \rho_j)}{T_{i+1} - T_{i-1}}
  \label{eq:Cv_fd}
\end{align}

\noindent\textbf{Axis-offset artefact note.} For a monatomic ideal gas,
$C_p = \frac{5}{2}R \approx 20.786\ \mathrm{J\,mol^{-1}\,K^{-1}}$ is
exactly constant. Any trajectory-integrated $C_p$ that varies from this value
by less than $\sim 10^{-8}\ \mathrm{J\,mol^{-1}\,K^{-1}}$ is numerical noise.
A chart renderer that re-scales the y-axis around these residuals \emph{must}
be rejected by the smoothness gate (\S\ref{sec:val_smooth}) and the chart
redrawn with a fixed axis spanning a physically meaningful range.

\subsection{Chart projections}

\begin{align}
  C_{\mathrm{sat}} &= \{(T_i,\, P_{\mathrm{sat},i})\}
  && \text{2D saturation curve} \label{eq:c_sat} \\[4pt]
  C_h             &= \{(T_i, P_j,\, h_{ij})\}
  && \text{3D enthalpy surface} \label{eq:c_h} \\[4pt]
  C_\rho          &= \{(T_i, P_j,\, \rho_{ij})\}
  && \text{3D density surface} \label{eq:c_rho} \\[4pt]
  \epsilon_h(T,P) &= h_{\mathrm{MD}}(T,P) - h_{\mathrm{ref}}(T,P)
  && \text{absolute residual} \label{eq:c_res} \\[4pt]
  \epsilon_{\mathrm{rel},h}(T,P) &=
    \frac{h_{\mathrm{MD}}(T,P) - h_{\mathrm{ref}}(T,P)}{h_{\mathrm{ref}}(T,P)}
  && \text{relative residual} \label{eq:c_relres}
\end{align}

% ============================================================================
\section{Molecular Dynamics Manifold}
\label{sec:md}
% ============================================================================

\subsection{Particle state}

\begin{equation}
  B_i = \bigl[\mathrm{id}_i,\; \xvec_i,\; \vvec_i,\; \avec_i,\;
               m_i,\; q_i,\; r_i,\; \mathbf{s}_i\bigr]
  \label{eq:particle}
\end{equation}
where $\avec_i = \Fvec_i / m_i$ is the acceleration.

\subsection{Force decomposition}

\begin{equation}
  \Fvec_i = \Fvec_{i,\mathrm{bonded}}
           + \Fvec_{i,\mathrm{near}}
           + \Fvec_{i,\mathrm{far}}
           + \Fvec_{i,\mathrm{flux}}
           + \Fvec_{i,\mathrm{event}}
  \label{eq:force_decomp}
\end{equation}

This decomposition is the bead communication spine: each term can be sourced
from a different module without restructuring the integrator.

\subsection{Time integration --- Velocity Verlet}

\begin{align}
  \xvec_{i,n+1} &= \xvec_{i,n} + \vvec_{i,n}\,\Delta t
    + \tfrac{1}{2}\avec_{i,n}\,\Delta t^2
  \label{eq:vv_x} \\[4pt]
  \vvec_{i,n+1} &= \vvec_{i,n}
    + \tfrac{1}{2}\bigl(\avec_{i,n} + \avec_{i,n+1}\bigr)\Delta t
  \label{eq:vv_v}
\end{align}

\subsection{Chart projections}

Trajectory:
\begin{equation}
  C_{\mathrm{traj}} = \{(\xvec_i(t),\, y_i(t),\, z_i(t))\}
  \label{eq:c_traj}
\end{equation}

Energy conservation residual:
\begin{equation}
  \epsilon_E(t) = \frac{E(t) - E(0)}{E(0)}
  \label{eq:c_edrift}
\end{equation}

Instantaneous kinetic temperature:
\begin{equation}
  T(t) = \frac{2}{3Nk_B}\sum_{i=1}^{N} \tfrac{1}{2}m_i\norm{\vvec_i}^2
  \label{eq:c_temp}
\end{equation}

Radial distribution function:
\begin{equation}
  g(r) = \frac{1}{4\pi r^2 \rho N}
  \left\langle \sum_{i} \sum_{j \neq i} \delta(r - r_{ij}) \right\rangle
  \label{eq:c_rdf}
\end{equation}

% ============================================================================
\section{Flux Field Manifold}
\label{sec:flux}
% ============================================================================

\subsection{Field types}

Scalar flux:
\begin{equation}
  \Phi_s = \Phi_s(\xvec, t, E)
  \label{eq:flux_scalar}
\end{equation}

Vector flux:
\begin{equation}
  \Phivec_v = \begin{pmatrix}\Phi_x \\ \Phi_y \\ \Phi_z\end{pmatrix}
  = \Phivec_v(\xvec, t, E)
  \label{eq:flux_vec}
\end{equation}

Tensor flux:
\begin{equation}
  \mathbf{\Phi}_T =
  \begin{pmatrix}
    \Phi_{xx} & \Phi_{xy} & \Phi_{xz} \\
    \Phi_{yx} & \Phi_{yy} & \Phi_{yz} \\
    \Phi_{zx} & \Phi_{zy} & \Phi_{zz}
  \end{pmatrix}
  \label{eq:flux_tensor}
\end{equation}

\subsection{Conservation form}

\begin{equation}
  \frac{\partial q}{\partial t} + \nabla \cdot \Phivec = S
  \label{eq:conservation}
\end{equation}

where $q$ is the conserved quantity density, $\Phivec$ is the flux, and $S$
is the source/sink term.

\subsection{Discrete flux divergence}

\begin{equation}
  \nabla \cdot \Phivec \approx
  \frac{\Phi_x(x+\Delta x) - \Phi_x(x-\Delta x)}{2\Delta x}
  + \frac{\Phi_y(y+\Delta y) - \Phi_y(y-\Delta y)}{2\Delta y}
  + \frac{\Phi_z(z+\Delta z) - \Phi_z(z-\Delta z)}{2\Delta z}
  \label{eq:div_discrete}
\end{equation}

\subsection{Primitive flux record}

\begin{lstlisting}[language=C++, caption={Primitive \texttt{.flux} record.}]
struct FluxSample {
    double   x, y, z;          // position (Ang)
    double   t;                 // time (fs)
    double   energy_eV;         // energy coordinate
    double   scalar;            // scalar field value
    double   vx, vy, vz;        // vector field components
    double   weight;            // statistical weight
    uint32_t field_type;        // flux type enum
    uint64_t source_event_id;   // parent event (0 = none)
};
\end{lstlisting}

\subsection{Chart projections}

2D slice at fixed $z = z_0$:
\begin{equation}
  C_{\Phi,\mathrm{2D}} = \{(x_i, y_i, \Phi_i) : z_i = z_0\}
  \label{eq:c_flux2d}
\end{equation}

3D vector field:
\begin{equation}
  C_{\Phi,\mathrm{vec}} = \{(\xvec_i, \Phivec_i)\}
  \label{eq:c_fluxvec}
\end{equation}

Divergence surface:
\begin{equation}
  C_{\nabla\cdot\Phi} = \{(x_i, y_i, (\nabla\cdot\Phivec)_i)\}
  \label{eq:c_div}
\end{equation}

% ============================================================================
\section{Radiation Transport Manifold}
\label{sec:radiation}
% ============================================================================

\subsection{Radiation state}

\begin{equation}
  \mathcal{R}_{\mathrm{rad}} = [\xvec,\; t,\; E,\; \Omvec,\;
                                  \Phi,\; D,\; \Sigma,\; \lambda,\; w]
  \label{eq:rad_state}
\end{equation}

\begin{center}
\begin{tabular}{lll}
  \toprule
  Symbol & Quantity & Unit \\
  \midrule
  $\Omvec$ & Direction unit vector & --- \\
  $\Phi$   & Particle flux         & m$^{-2}$\,s$^{-1}$ \\
  $D$      & Dose / deposited energy & J/kg (Gy) \\
  $\Sigma$ & Macroscopic cross section & m$^{-1}$ \\
  $\lambda = 1/\Sigma$ & Attenuation length & m \\
  $w$      & Particle weight & --- \\
  \bottomrule
\end{tabular}
\end{center}

\subsection{Attenuation law}

\begin{equation}
  I(x) = I_0\, e^{-\mu x}, \qquad
  T(x) = \frac{I(x)}{I_0} = e^{-\mu x}
  \label{eq:attenuation}
\end{equation}

\subsection{Flux definitions}

Integrated flux:
\begin{equation}
  \Phi = \frac{dN}{dA\,dt}
  \label{eq:flux_def}
\end{equation}

Energy-resolved flux:
\begin{equation}
  \Phi(E) = \frac{dN}{dA\,dt\,dE}
  \label{eq:flux_energy}
\end{equation}

\subsection{Dose proxy}

\begin{equation}
  D = \frac{E_{\mathrm{dep}}}{m}, \qquad
  \dot{D} = \frac{dD}{dt}
  \label{eq:dose}
\end{equation}

\subsection{Chart projections}

Energy spectrum, dose map, flux field, and attenuation curve:
\begin{align}
  C_E         &= \{(E_i, N_i)\}
  && \text{energy spectrum} \label{eq:c_espect} \\[4pt]
  C_D         &= \{(x_i, y_i, D_i) : z = z_0\}
  && \text{dose map} \label{eq:c_dose} \\[4pt]
  C_\Phi      &= \{(\xvec_i, \Phi_i)\}
  && \text{flux field} \label{eq:c_flux} \\[4pt]
  C_T         &= \{(x_i,\; e^{-\mu x_i})\}
  && \text{attenuation curve} \label{eq:c_atten}
\end{align}

% ============================================================================
\section{Excitation Dynamics Manifold}
\label{sec:excitation}
% ============================================================================

\subsection{Excitation state}

\begin{equation}
  \mathcal{X}_i = [\mathrm{id}_i,\; t_i,\; \xvec_i,\;
                    E_{0,i},\; E_{1,i},\; \Delta E_i,\;
                    \sigma_i,\; p_i,\; \tau_i]
  \label{eq:excit_state}
\end{equation}
where $\Delta E_i = E_{1,i} - E_{0,i}$, $\sigma_i$ is the transition
cross-section, $p_i$ is the transition probability, and $\tau_i$ is the
state lifetime.

\subsection{Population dynamics}

State population vector with normalisation:
\begin{equation}
  \mathbf{p}(t) = [p_1(t),\; p_2(t),\; \ldots,\; p_n(t)],
  \qquad
  \sum_{k=1}^{n} p_k(t) = 1
  \label{eq:population}
\end{equation}

Transition rate matrix $K$ and population equation:
\begin{equation}
  \frac{d\mathbf{p}}{dt} = K\,\mathbf{p},
  \qquad
  \mathbf{p}_{n+1} = \mathbf{p}_n + \Delta t\, K\,\mathbf{p}_n
  \label{eq:rate_eq}
\end{equation}

\subsection{Excitation event struct}

\begin{lstlisting}[language=C++, caption={Excitation event in \texttt{.dynx}.}]
struct ExcitationEvent {
    uint64_t event_id;
    uint64_t particle_id;
    double   t;                 // time (fs)
    double   x, y, z;          // position (Ang)
    double   E_initial_eV;
    double   E_final_eV;
    double   delta_E_eV;
    int      state_initial;
    int      state_final;
    uint32_t cause;             // collision | photon_abs | decay | thermal | field
};
\end{lstlisting}

\subsection{Chart projections}

\begin{align}
  C_\mathcal{X}  &= \{(t_i,\; \Delta E_i)\}
  && \text{transition timeline} \label{eq:c_xtl} \\[4pt]
  C_p            &= \{(t_i,\; p_k(t_i))\}
  && \text{state population} \label{eq:c_pop} \\[4pt]
  C_{\mathrm{emit}} &= \{(\Delta E_i,\; N_i)\}
  && \text{emission spectrum} \label{eq:c_emit}
\end{align}

% ============================================================================
\section{Effective Gluon / Color-Field Proxy Manifold}
\label{sec:gluon}
% ============================================================================

\begin{quote}
  \itshape
  \textbf{Scope disclaimer.} The model in this section is an
  \textbf{effective proxy}, not a first-principles QCD calculation.
  It does not solve the QCD path integral, does not implement lattice
  gauge theory, and makes no claim to reproduce the full non-Abelian
  $\mathrm{SU}(3)$ dynamics. It is a reduced representation suitable
  for event-level bookkeeping in particle cascade simulations and for
  visualisation of confinement-like phenomena in educational or
  exploratory contexts. The field \texttt{full\_qcd = false} must always
  be set in any \texttt{.xchart} or \texttt{.dynx} output from this module.
  The universe already contains enough fraud without compiling more.
\end{quote}

\subsection{Effective color-state vector}

The six-component color-state vector encodes the three color and three
anti-color occupancy proxies:
\begin{equation}
  \mathbf{G}_i =
  \bigl[g_r,\; g_g,\; g_b,\; \bar{g}_r,\; \bar{g}_g,\; \bar{g}_b\bigr]
  \in [0,1]^6
  \label{eq:color_vec}
\end{equation}

or, using the reduced scalar for confinement dynamics, the color-excitation
state:
\begin{equation}
  \chi_c(t) \in [0,1],
  \qquad
  \chi_c(t) = 0 \implies \text{color neutral (confined)}
  \label{eq:chi_c}
\end{equation}

\subsection{Color-neutrality residual}

The residual from exact color neutrality is:
\begin{equation}
  \eta_c = \left\| \sum_{i \in \mathrm{hadron}} \mathbf{G}_i \right\|,
  \qquad
  \eta_c \leq \epsilon_c \implies \text{color-neutral state}
  \label{eq:color_residual}
\end{equation}

\subsection{Confinement proxy energy}

A phenomenological confinement energy proxy:
\begin{equation}
  E_{\mathrm{conf}}(r) = \kappa\, r + \frac{\alpha_s}{r},
  \qquad
  \kappa \approx 0.18\ \mathrm{GeV}^2,\quad
  \alpha_s \approx 0.1\text{--}0.3
  \label{eq:confinement}
\end{equation}

This is the Cornell potential. It is an effective model. It is not derived
from the VSEPR-SIM force field and should not be connected to the MD
integrator without explicit demarcation.

\subsection{Gluon event struct}

\begin{lstlisting}[language=C++, caption={Effective gluon excitation event.}]
struct GluonExcitationEvent {
    uint64_t event_id;
    uint64_t parent_particle_id;
    double   t;                       // time (s)
    double   x, y, z;                 // position (m)
    double   energy_GeV;              // effective energy
    double   confinement_proxy;       // Cornell potential value
    double   color_state[6];          // G_i vector
    uint32_t interaction_type;        // emission | absorption | annihilation
    uint64_t state_hash;              // provenance
};
\end{lstlisting}

\subsection{Chart projections}

\begin{align}
  C_{\chi}   &= \{(t_i,\; \chi_{c,i})\}
  && \text{color-excitation timeline} \label{eq:c_chi} \\[4pt]
  C_{\mathrm{conf}} &= \{(r_i,\; E_{\mathrm{conf},i})\}
  && \text{confinement curve} \label{eq:c_conf} \\[4pt]
  C_G        &= \{(t_i,\; \mathbf{G}_i)\}
  && \text{color-state history} \label{eq:c_gluon}
\end{align}

% ============================================================================
\section{Dynamic Event Stream --- \texttt{.dynx}}
\label{sec:dynx}
% ============================================================================

\subsection{General event object}

\begin{equation}
  D_i = [\mathrm{id}_i,\; \mathrm{type}_i,\; t_i,\; \xvec_i,\;
          \avec_i,\; \mathbf{b}_i,\; \mathbf{q}_i,\; H_i]
  \label{eq:dynx_event}
\end{equation}

where $\avec_i$ is the primary vector payload, $\mathbf{b}_i$ is the
secondary vector payload, and $\mathbf{q}_i$ is the scalar payload vector.

\subsection{Bond-force vector event}

\begin{equation}
  D_{\mathrm{bond}} = [\mathrm{id},\; t,\; \xvec_a,\; \xvec_b,\;
                        \Fvec_{ab},\; \ell,\; \ell_0,\; \epsilon_b]
  \label{eq:bond_event}
\end{equation}

Bond length and strain:
\begin{align}
  \ell      &= \norm{\xvec_b - \xvec_a} \label{eq:bond_len} \\
  \epsilon_b &= \frac{\ell - \ell_0}{\ell_0} \label{eq:bond_strain}
\end{align}

Bond force (harmonic):
\begin{equation}
  \Fvec_{ab} = -k_b\,(\ell - \ell_0)\,
               \frac{\xvec_b - \xvec_a}{\ell}
  \label{eq:bond_force}
\end{equation}

Break criteria:
\begin{equation}
  |\epsilon_b| > \epsilon_{\mathrm{crit}}
  \quad\text{or}\quad
  \norm{\Fvec_{ab}} > F_{\mathrm{crit}}
  \label{eq:bond_break}
\end{equation}

\subsection{\texttt{.dynx} data contract}

\begin{center}
\begin{tabular}{ll}
  \toprule
  Data-in field & Meaning \\
  \midrule
  \texttt{event\_id}, \texttt{event\_type} & Event identity \\
  \texttt{time\_s}                         & Event time \\
  \texttt{particle\_a}, \texttt{particle\_b} & Participant IDs \\
  \texttt{x, y, z}                         & Event position \\
  \texttt{vx, vy, vz}                      & Velocity at event \\
  \texttt{Fx, Fy, Fz}                      & Force at event \\
  \texttt{energy}                          & Event energy \\
  \texttt{state\_initial}, \texttt{state\_final} & State transition \\
  \texttt{hash}                            & Event provenance hash \\
  \bottomrule
\end{tabular}
\end{center}

% ============================================================================
\section{File Format Specifications}
\label{sec:formats}
% ============================================================================

\subsection{\texttt{.chart} --- lightweight plot recipe}

A \texttt{.chart} file defines the tuple:
\begin{equation}
  C = (S,\; \Proj,\; \mathcal{A},\; \Val)
  \label{eq:chart_tuple}
\end{equation}

where $S$ is the source dataset path, $\Proj$ is the projection (column
mapping), $\mathcal{A}$ is the axis metadata set, and $\Val$ is the
validation rule set.

Axis definition:
\begin{equation}
  \mathcal{A}_k = (\mathrm{name}_k,\; \mathrm{unit}_k,\; f_k)
  \label{eq:axis}
\end{equation}

\begin{lstlisting}[caption={\texttt{.chart} format example.}]
[chart]
version    = "0.2"
kind       = "line"
title      = "Cp(T) - Hydrogen Atom"

[source]
path       = "h_atom_steam.tsv"
x_column   = "T_K"
y_column   = "Cp_J_molK"

[axes]
x_label    = "Temperature"
x_unit     = "K"
x_min      = 100.0
x_max      = 6000.0
y_label    = "Cp"
y_unit     = "J/(mol K)"
y_min      = 20.0          # fixed --- prevents offset artifact
y_max      = 22.0

[validation]
y_finite        = true
y_nonnegative   = true
smoothness_tol  = 0.05
axis_offset_check = true   # rejects plots where offset > 1e-6 * range
\end{lstlisting}

\subsection{\texttt{.xchart} --- rich scientific visualization object}

An \texttt{.xchart} file defines:
\begin{equation}
  \mathcal{X} = (\Tab,\; \mathcal{P},\; \mathcal{M},\; \Val,\; H)
  \label{eq:xchart}
\end{equation}

where $\mathcal{P} = \{\Proj_1, \Proj_2, \ldots, \Proj_k\}$ is the set of
chart projections, $\mathcal{M}$ is the metadata block, and $H$ is the
provenance hash.

\begin{lstlisting}[caption={\texttt{.xchart} --- radiation flux surface.}]
[xchart]
version = "0.2"
kind    = "radiation_flux_surface"
title   = "Neutron Flux Field Around Source"

[datasets.primary]
format  = "inline_table"
columns = ["x_m","y_m","z_m","time_s","energy_MeV",
           "flux_m2_s","dose_Gy_s"]

[plot.surface]
type = "3d_surface"
x    = "radius_m"
y    = "energy_MeV"
z    = "flux_m2_s"

[physics]
field           = "neutron_flux"
transport_model = "effective_monte_carlo"
source          = "decay_event_stream"
units_checked   = true

[validation]
nonneg_flux    = true
energy_bins_ok = true
source_hash    = "PLACEHOLDER"
\end{lstlisting}

\begin{lstlisting}[caption={\texttt{.xchart} --- effective gluon excitation timeline.}]
[xchart]
version = "0.2"
kind    = "subatomic_excitation_timeline"
title   = "Effective Gluon Excitation Events"

[datasets.events]
format  = "inline_table"
columns = ["event_id","particle_id","time_s","energy_GeV",
           "confinement_proxy","color_r","color_g","color_b",
           "anticolor_r","anticolor_g","anticolor_b"]

[plot.timeline]
type  = "event_timeline"
x     = "time_s"
y     = "energy_GeV"
group = "particle_id"

[plot.phase]
type = "phase_map"
x    = "time_s"
y    = "confinement_proxy"
z    = "energy_GeV"

[physics]
model_class = "effective_subatomic_proxy"
full_qcd    = false          # mandatory --- this is NOT a QCD solver
\end{lstlisting}

\subsection{\texttt{.flux} --- scalar/vector/tensor field state}

The \texttt{.flux} file is a binary or text stream of
\texttt{FluxSample} records (see \S\ref{sec:flux}). Text form uses
tab-separated columns in the order defined by the file header:

\begin{lstlisting}[caption={\texttt{.flux} text header.}]
# flux_version=1
# field_type=neutron_flux
# units=x:m,t:s,energy:eV,phi:m-2.s-1
x_m  y_m  z_m  time_s  energy_eV  phi_scalar  phi_x  phi_y  phi_z  weight  source_id
\end{lstlisting}

\subsection{\texttt{.x} --- universal session bundle}

The bundle is:
\begin{equation}
  \Bund = \{S_{\mathrm{sim}},\; \mathcal{T}_{\mathrm{state}},\;
             D_{\mathrm{dyn}},\; \mathcal{F}_{\mathrm{flux}},\;
             C_{\mathrm{chart}},\; R_{\mathrm{report}},\; H_\Bund\}
  \label{eq:bundle}
\end{equation}

Bundle hash:
\begin{equation}
  H_\Bund = \Hash(H_1 \,\|\, H_2 \,\|\, \cdots \,\|\, H_n)
  \label{eq:bundle_hash}
\end{equation}

\begin{center}
\begin{tabular}{lll}
  \toprule
  Symbol & Contents & Extensions \\
  \midrule
  $S_{\mathrm{sim}}$        & Simulation scripts     & \texttt{.vsim} \\
  $\mathcal{T}_{\mathrm{state}}$ & Particle states   & \texttt{.xyz}, \texttt{.xyza}, \texttt{.xyzf}, \texttt{.xyzc} \\
  $D_{\mathrm{dyn}}$        & Dynamic event streams  & \texttt{.dynx} \\
  $\mathcal{F}_{\mathrm{flux}}$ & Field data         & \texttt{.flux} \\
  $C_{\mathrm{chart}}$      & Charts / visualizations & \texttt{.chart}, \texttt{.xchart} \\
  $R_{\mathrm{report}}$     & Reports and logs       & \texttt{.json}, \texttt{.tsv}, \texttt{.md} \\
  \bottomrule
\end{tabular}
\end{center}

% ============================================================================
\section{Qt Desktop Stack}
\label{sec:qt}
% ============================================================================

\subsection{View modes}

The Qt desktop layer treats charts as scientific view modes, not plots.
Each view kind maps to a specific projection type from \S\ref{sec:projection}.

\begin{center}
\begin{tabular}{lll}
  \toprule
  Qt view class & Projection type & Primary input \\
  \midrule
  \texttt{Chart2DView}           & $C_{\mathrm{2D}}$         & \texttt{.chart}, \texttt{.xchart} \\
  \texttt{Chart3DSurfaceView}    & $C_{\mathrm{3D}}$         & \texttt{.xchart} \\
  \texttt{FieldSliceView}        & $C_{\Phi,\mathrm{2D}}$    & \texttt{.flux}, \texttt{.xchart} \\
  \texttt{FluxVectorView}        & $C_{\Phi,\mathrm{vec}}$   & \texttt{.flux} \\
  \texttt{RadiationMapView}      & $C_D$, $C_\Phi$           & \texttt{.xchart}, \texttt{.dynx} \\
  \texttt{ExcitationTimelineView}& $C_\mathcal{X}$           & \texttt{.dynx} \\
  \texttt{SubatomicEventView}    & $C_\chi$, $C_G$           & \texttt{.dynx} \\
  \texttt{BatchComparisonView}   & $\{\Proj_j\}_j$           & \texttt{.x} bundle \\
  \bottomrule
\end{tabular}
\end{center}

\subsection{Desktop opening behaviour}

\begin{lstlisting}[caption={Qt desktop launch examples.}]
vsepr open steam_surface.xchart        # -> Chart3DSurfaceView
vsepr open radiation_flux.xchart       # -> RadiationMapView
vsepr open gluon_excitation.dynx       # -> SubatomicEventView + FluxVectorView
vsepr open reactor_batch.x             # -> BatchComparisonView dashboard
\end{lstlisting}

\subsection{Visualisation families}

\textbf{Thermodynamic charts.} For steam tables and MD-derived properties.
Views: 2D saturation curve, 3D property surface, phase map, residual error
map, lookup interpolation heatmap.

\textbf{Radiation charts.} For isotope, reactor, shielding, and particle
transport. Variables: $\Phi(r,t,E) = dN/(dA\,dt\,dE)$. Views: flux vs.\
radius, dose vs.\ shielding thickness, energy spectrum, attenuation curve,
3D radiation field.

\textbf{Flux fields.} For thermal, mass, charge, neutron, momentum, and
reaction flux. Connects directly to \texttt{.flux}. Views: slice, divergence
surface, vector arrows.

\textbf{Excitation dynamics.} For atoms, molecules, radiation events, and
reduced quantum states. Views: excitation energy vs.\ time, transition
probability map, emission spectrum, state population histogram, decay cascade
timeline.

\textbf{Effective gluon / subatomic events.} Marked \texttt{full\_qcd =
false}. Views: color-state transition timeline, effective confinement energy,
subatomic excitation graph, collision event tree, particle cascade map, field
intensity surface.

% ============================================================================
\section{Validation Mathematics}
\label{sec:validation}
% ============================================================================

\subsection{Unit consistency}

Every plotted variable must satisfy dimensional agreement with its declared
axis:
\begin{equation}
  \dim(f_k) = \dim(\mathcal{A}_k) \quad \forall\, k
  \label{eq:val_dim}
\end{equation}
Reject if $\dim(f_k) \neq \dim(\mathcal{A}_k)$ for any $k$.

\subsection{Finite-value gate}

\begin{equation}
  \forall\, i,k : \mathrm{isfinite}\bigl(f_k(\Rec_i)\bigr) = \mathrm{true}
  \label{eq:val_finite}
\end{equation}
Reject if any sample maps to $\{\mathrm{NaN},\, +\infty,\, -\infty\}$.

\subsection{Non-negativity gate}

For density, absolute temperature, flux magnitude, dose, and mass:
\begin{equation}
  q_i \geq 0 \quad \forall\, i
  \label{eq:val_nonneg}
\end{equation}

\subsection{Conservation residual}

For conserved density $q$ governed by \eqref{eq:conservation}:
\begin{equation}
  \mathcal{R}_q = \frac{\partial q}{\partial t} + \nabla\cdot\Phivec - S
  \label{eq:val_conservation}
\end{equation}
Accept if $|\mathcal{R}_q| < \epsilon_q$.

\subsection{Smoothness gate}
\label{sec:val_smooth}

Second-difference smoothness check:
\begin{equation}
  \Delta^2 y_i = y_{i+1} - 2y_i + y_{i-1}
  \label{eq:val_smooth}
\end{equation}
Accept if $|\Delta^2 y_i| < \epsilon_{\mathrm{smooth}}$ unless
$\mathrm{boundary}_i = \mathrm{true}$ (phase boundary).

\noindent\textbf{Application to the Cp axis artifact.} For a monatomic
ideal-gas $C_p$ trace, $|\Delta^2 C_p| \ll 10^{-6}\ \mathrm{J\,mol^{-1}\,K^{-1}}$.
A chart renderer that produces an axis offset greater than
$10^{-6} \times y_{\mathrm{range}}$ on such a trace must fail this gate and
re-render with fixed axis bounds.

\subsection{Segment discontinuity gate}

At a segment boundary $x_b$ between piecewise model regions:
\begin{equation}
  \Delta y_b = \left| y^+(x_b) - y^-(x_b) \right|
  \label{eq:val_disc}
\end{equation}
Accept if $\Delta y_b < \epsilon_b$ or $\mathrm{phase\_boundary}_b =
\mathrm{true}$.

\noindent This gate would have caught the ethane/propane plot discontinuities
at segment boundaries in the previous implementation.

\subsection{Axis-offset rejection gate}

Let $y_{\mathrm{range}} = y_{\max} - y_{\min}$ over the plotted data. Define
the axis offset as the mantissa displacement introduced by the renderer:
\begin{equation}
  \delta_{\mathrm{offset}} = \bigl|y_{\mathrm{axis,min}} - y_{\min}\bigr|
  \label{eq:val_offset}
\end{equation}
Reject (flag as artifact) if:
\begin{equation}
  \frac{\delta_{\mathrm{offset}}}{y_{\mathrm{range}}} < \epsilon_{\mathrm{offset}}
  \quad \text{and} \quad
  \delta_{\mathrm{offset}} > 0
  \label{eq:val_offset_gate}
\end{equation}
Default $\epsilon_{\mathrm{offset}} = 10^{-4}$. A non-zero offset on a
nearly-flat curve indicates the renderer is auto-scaling around noise.

% ============================================================================
\section{Universal Pipeline and Implementation Contract}
\label{sec:contract}
% ============================================================================

\subsection{Universal pipeline}

Every new module must implement the four-stage pipeline:

\begin{equation}
  \mathcal{I} \;\xrightarrow{K}\; \Tab
                \;\xrightarrow{\Proj_j}\; C_j
                \;\xrightarrow{\Val_j}\; \{\mathrm{pass},\,\mathrm{fail},\,\epsilon\}
                \;\xrightarrow{\mathcal{E}}\; \{\texttt{.chart},\,\texttt{.xchart},\,\ldots\}
  \label{eq:pipeline}
\end{equation}

\begin{description}[leftmargin=2em]
  \item[Input $\mathcal{I}$] Raw simulation or measurement records
    $\{\Rec_i^{\mathrm{raw}}\}$.
  \item[Transform $K$] Normalise to the universal record schema
    \eqref{eq:record}: unit conversion, column mapping, hash computation.
  \item[Project $\Proj_j$] Apply one or more projection operators to produce
    chart tables $C_j$.
  \item[Validate $\Val_j$] Apply all applicable gates from \S\ref{sec:validation}.
    A failing gate produces a diagnostic, not a silent bad chart.
  \item[Export $\mathcal{E}$] Write \texttt{.chart}, \texttt{.xchart},
    \texttt{PNG}, \texttt{SVG}, \texttt{TSV}, or bundle into \texttt{.x}.
\end{description}

In plain structure:

\begin{lstlisting}[caption={Universal pipeline pseudocode.}]
raw_data
  -> normalize_to_universal_table(K)
  -> for each projection Proj_j:
       C_j = project(table, Proj_j)
       result = validate(C_j, rules_j)
       if result.pass:
           export(C_j, fmt)
       else:
           emit_diagnostic(result)
\end{lstlisting}

\subsection{Variable dictionary}

The following variable names are canonical across all modules and file formats:

\begin{center}
\begin{tabular}{ll}
  \toprule
  Name & Meaning \\
  \midrule
  \texttt{id}                     & Persistent sample or particle identifier \\
  \texttt{x, y, z}                & Position (\AA\ for MD; m for radiation) \\
  \texttt{t}                      & Time \\
  \texttt{T}                      & Temperature (K) \\
  \texttt{P}                      & Pressure (Pa) \\
  \texttt{rho}                    & Density (kg/m$^3$) \\
  \texttt{u}                      & Specific internal energy (J/kg) \\
  \texttt{h}                      & Specific enthalpy (J/kg) \\
  \texttt{s}                      & Specific entropy (J/(kg\,K)) \\
  \texttt{Cp, Cv}                 & Heat capacities (J/(mol\,K)) \\
  \texttt{mu}                     & Dynamic viscosity (Pa\,s) \\
  \texttt{k}                      & Thermal conductivity (W/(m\,K)) \\
  \texttt{E}                      & Energy (eV for MD; MeV/GeV for particle) \\
  \texttt{vx, vy, vz}             & Velocity (\AA/fs) \\
  \texttt{Fx, Fy, Fz}             & Force (eV/\AA) \\
  \texttt{Phi}                    & Scalar flux \\
  \texttt{Phi\_x, Phi\_y, Phi\_z} & Vector flux components \\
  \texttt{D}                      & Dose (Gy) \\
  \texttt{Omega}                  & Radiation direction unit vector \\
  \texttt{Sigma}                  & Macroscopic cross section (m$^{-1}$) \\
  \texttt{state\_i, state\_j}     & Initial and final excitation state indices \\
  \texttt{delta\_E}               & Transition energy \\
  \texttt{c}                      & Color-state proxy vector $\mathbf{G}_i$ \\
  \texttt{eta\_c}                 & Color-neutrality residual \\
  \texttt{E\_conf}                & Confinement proxy energy \\
  \texttt{w}                      & Sample weight \\
  \texttt{epsilon}                & Validation residual \\
  \texttt{H}                      & Provenance hash \\
  \bottomrule
\end{tabular}
\end{center}

\subsection{Implementation order}

\begin{enumerate}
  \item Implement \texttt{.chart} format parser and writer.
  \item Implement \texttt{.xchart} format with embedded table support.
  \item Implement \texttt{.flux} binary and text record I/O.
  \item Connect \texttt{.xchart} to \texttt{.flux} (field visualization).
  \item Connect \texttt{.xchart} to \texttt{.dynx} (event timeline visualization).
  \item Add Qt \texttt{Chart2DView} and \texttt{Chart3DSurfaceView}.
  \item Add Qt \texttt{FluxVectorView} and \texttt{FieldSliceView}.
  \item Add Qt \texttt{ExcitationTimelineView}.
  \item Add Qt \texttt{RadiationMapView} and \texttt{SubatomicEventView}.
  \item Add Qt \texttt{BatchComparisonView}.
  \item Add all validation gates (\S\ref{sec:validation}).
  \item Bundle all outputs through \texttt{.x}.
\end{enumerate}

\subsection{Data-in / data-out matrix}

\begin{center}
\begin{tabular}{p{3.2cm}p{4.8cm}p{4.8cm}}
  \toprule
  Module & Data-in & Data-out \\
  \midrule
  Steam MD table
    & $T, P, \rho, u, h, s, C_p, C_v, \mu, k$
    & \texttt{.xchart}, lookup cache, phase maps \\
  Molecular dynamics
    & $\xvec, \vvec, m, q, \Fvec$, forcefield
    & \texttt{.xyzf}, \texttt{.dynx}, \texttt{.xchart} \\
  Flux field
    & $\xvec, t, E, \Phi, \Phivec$
    & \texttt{.flux}, vector/slice charts \\
  Radiation transport
    & particle type, $E, \Omvec, \Sigma, \Phi, D$
    & \texttt{.dynx}, \texttt{.flux}, dose maps \\
  Excitation dynamics
    & $E_i, E_j, \Delta E, p, \tau$
    & timelines, emission spectra \\
  Effective gluon proxy
    & $\mathbf{G}, E_{\mathrm{conf}}, \chi_c$
    & color-state maps, event trees \\
  Chart engine
    & source columns, axes, units, rules
    & PNG/SVG/Qt view \\
  XChart engine
    & embedded datasets, projections, hashes
    & multi-view scientific artifact \\
  Bundle \texttt{.x}
    & all artifact types
    & replayable session archive \\
  \bottomrule
\end{tabular}
\end{center}

\subsection{Architecture rule}

Do not make gluons, radiation, flux, and steam tables separate isolated
plotting systems. They are all projections of the universal state vector
$\Qvec$ \eqref{eq:Qvec} onto appropriate subspaces:

\begin{equation}
  \mathcal{D} = \{X, Y, Z, T, E, S, V, M, H\}
  \label{eq:arch_rule}
\end{equation}

where $X,Y,Z$ are position; $T$ is time or temperature; $E$ is energy; $S$
are state variables; $V$ are vector field components; $M$ is model metadata;
and $H$ is the provenance hash. Then \texttt{.xchart} maps those fields into
views. One spine. All output types fall from it.

\vfill
\begin{center}
  \rule{0.4\textwidth}{0.4pt}\\[0.5em]
  \textit{WO-CHART-FIELD-B --- VSEPR-SIM Theoretical Division.
  One abstraction. All the physics. Someone alert the academy before
  they ruin it with a subcommittee.}
\end{center}

\end{document}
